\documentclass[11pt,a4paper]{article}

% ----------------------------------
% Packages
% ----------------------------------
\usepackage[utf8]{inputenc}
\usepackage[T1]{fontenc}
\usepackage[a4paper,top=2.1cm,bottom=2.1cm,left=2.25cm,right=2.25cm]{geometry}
\usepackage{xcolor}
\usepackage{hyperref}
\usepackage{enumitem}
\usepackage{titlesec}
\usepackage{fancyhdr}
\usepackage{listings}
\usepackage{listingsutf8}
\usepackage{bookmark}
\usepackage{lmodern}
\usepackage{microtype}
\usepackage{upquote}

% ----------------------------------
% Document information
% ----------------------------------
\newcommand{\DocumentType}{Lab Writeups}
\newcommand{\VulnerabilityTitle}{API Testing}
\newcommand{\SourceName}{PortSwigger Web Security Academy}
\newcommand{\AuthorName}{Youssef Reda}

% ----------------------------------
% Colors
% ----------------------------------
\definecolor{primary}{RGB}{25,55,95}
\definecolor{secondary}{RGB}{75,95,120}
\definecolor{linkblue}{RGB}{35,85,155}
\definecolor{codegray}{RGB}{95,95,95}
\definecolor{codebg}{RGB}{247,248,250}
\definecolor{codeframe}{RGB}{210,216,224}

% ----------------------------------
% Code listings
% ----------------------------------
\lstset{
    inputencoding=utf8,
    basicstyle=\ttfamily\small,
    backgroundcolor=\color{codebg},
    frame=single,
    rulecolor=\color{codeframe},
    framesep=7pt,
    xleftmargin=4pt,
    xrightmargin=4pt,
    aboveskip=8pt,
    belowskip=8pt,
    breaklines=true,
    breakatwhitespace=false,
    columns=fullflexible,
    keepspaces=true,
    showstringspaces=false,
    tabsize=2
}

\lstdefinelanguage{JavaScript}{
    language=Java,
    basicstyle=\ttfamily\small,
    numbers=none,
    showstringspaces=false,
    morekeywords={const,let,var,async,await,fetch,return,import,export,from},
    keywordstyle=\color{primary}\bfseries,
    stringstyle=\color{secondary},
    morestring=[b]',
    morestring=[b]"
}

\lstdefinelanguage{JSON}{
    basicstyle=\ttfamily\small,
    numbers=none,
    showstringspaces=false,
    morestring=[b]",
    stringstyle=\color{secondary}
}

\lstdefinelanguage{HTTP}{
    basicstyle=\ttfamily\small,
    numbers=none,
    showstringspaces=false,
    morestring=[b]",
    stringstyle=\color{secondary},
    morecomment=[l]{\#}
}

% ----------------------------------
% Hyperlinks and PDF metadata
% ----------------------------------
\hypersetup{
    colorlinks=true,
    linkcolor=primary,
    urlcolor=linkblue,
    citecolor=primary,
    pdftitle={\VulnerabilityTitle},
    pdfauthor={\AuthorName}
}

\urlstyle{same}

% ----------------------------------
% Paragraphs, lists, and headings
% ----------------------------------
\setlength{\parindent}{0pt}
\setlength{\parskip}{4pt}

\setlist[itemize]{
    leftmargin=1.5em,
    itemsep=2pt,
    topsep=3pt
}

\setlist[enumerate]{
    leftmargin=1.8em,
    itemsep=2pt,
    topsep=3pt
}

\titleformat{\section}
  {\Large\bfseries\color{primary}}{}{0pt}{}

\titlespacing*{\section}{0pt}{16pt}{6pt}

% ----------------------------------
% Unnumbered headings in Table of Contents
% ----------------------------------
\newcommand{\nsection}[1]{%
  \section*{#1}%
  \phantomsection%
  \addcontentsline{toc}{section}{#1}%
}

\setcounter{tocdepth}{1}

% ----------------------------------
% Header and footer
% ----------------------------------
\pagestyle{fancy}
\fancyhf{}

\fancyhead[L]{\small\textcolor{secondary}{Web Security Notes}}
\fancyhead[R]{\small\textit{API Labs}}

\fancyfoot[C]{\thepage}

\renewcommand{\headrulewidth}{0.4pt}
\renewcommand{\footrulewidth}{0pt}

\setlength{\headheight}{14pt}

% ----------------------------------
% Document
% ----------------------------------
\begin{document}

% ----------------------------------
% Cover page
% ----------------------------------
\begin{titlepage}
    \thispagestyle{empty}
    \centering

    \vspace*{1.8cm}
    {\Large\color{secondary}\bfseries \DocumentType\par}

    \vspace{1.2cm}
    {\Huge\bfseries\color{primary}\VulnerabilityTitle\par}

    \vspace{0.8cm}
    {\large PortSwigger Web Security Academy — Lab Solutions\par}

    \vfill
    \begin{flushleft}
        \vspace{0.6cm}
        \textbf{Prepared by}\par
        \AuthorName
    \end{flushleft}
\end{titlepage}

% ----------------------------------
% Table of Contents
% ----------------------------------
\pagenumbering{roman}
\tableofcontents

\clearpage
\pagenumbering{arabic}

% ----------------------------------
% Lab Writeups
% ----------------------------------

\nsection{Lab: Exploiting an API Endpoint Using Documentation}

\begin{enumerate}
    \item Log in to the application using the credentials \texttt{wiener:peter} and update your email address.
    \item In \textbf{Proxy > HTTP history}, locate the \texttt{PATCH /api/user/wiener} request. Right-click it and select \textbf{Send to Repeater}.
    \item In Repeater, send the request and observe that it retrieves credentials for the user \texttt{wiener}.
    \item Remove \texttt{/wiener} from the request path:
          \begin{lstlisting}[language=HTTP]
PATCH /api/user
    \end{lstlisting}
          Send the request and observe the error indicating that a user identifier is required.
    \item Remove \texttt{/user} from the path:
          \begin{lstlisting}[language=HTTP]
GET /api
    \end{lstlisting}
          Send the request and observe that it returns API documentation.
    \item Right-click the response and select \textbf{Show response in browser}. Copy the generated URL.
    \item Paste the URL into Burp's browser to access the interactive API documentation.
    \item Locate the \texttt{DELETE} endpoint, enter \texttt{carlos} as the user identifier, and click \textbf{Send request} to delete the user and solve the lab.
\end{enumerate}

\newpage
\nsection{Lab: Finding and Exploiting an Unused API Endpoint}

\begin{enumerate}
    \item Access the lab in Burp's browser and click on any product.
    \item In \textbf{Proxy > HTTP history}, locate the API request associated with the product, for example:
          \begin{lstlisting}[language=HTTP]
GET /api/products/3/price
    \end{lstlisting}
    \item Right-click the request and select \textbf{Send to Repeater}.
    \item Change the HTTP method from \texttt{GET} to \texttt{OPTIONS} and send the request. Observe that the response indicates that \texttt{GET} and \texttt{PATCH} methods are allowed.
    \item Change the HTTP method to \texttt{PATCH} and send the request. Observe the \texttt{Unauthorized} response, indicating that authentication may be required.
    \item Log in to the application using the credentials \texttt{wiener:peter}.
    \item Open the \textbf{Lightweight "l33t" Leather Jacket} product.
    \item In \textbf{Proxy > HTTP history}, locate the following request and send it to Repeater:
          \begin{lstlisting}[language=HTTP]
GET /api/products/1/price
    \end{lstlisting}
    \item Change the HTTP method to \texttt{PATCH} and send the request. Observe the error indicating that the expected Content-Type is \texttt{application/json}.
    \item Add the following header:
          \begin{lstlisting}[language=HTTP]
Content-Type: application/json
    \end{lstlisting}
    \item Add an empty JSON object as the request body:
          \begin{lstlisting}[language=JSON]
{}
    \end{lstlisting}
          Send the request and observe the error indicating that the \texttt{price} parameter is required.
    \item Add the \texttt{price} parameter with a value of \texttt{0}:
          \begin{lstlisting}[language=JSON]
{
    "price": 0
}
    \end{lstlisting}
    \item Send the request. In Burp's browser, reload the leather jacket product page and verify that its price is now \$0.00.
    \item Add the leather jacket to your basket.
    \item Open the basket and click \textbf{Place order} to solve the lab.
\end{enumerate}

\newpage
\nsection{Lab: Exploiting a Mass Assignment Vulnerability}

\begin{enumerate}
    \item Log in to the application using the credentials \texttt{wiener:peter}.
    \item Open the \textbf{Lightweight "l33t" Leather Jacket} product and add it to your basket.
    \item Open the basket and click \textbf{Place order}. Observe that the account does not have enough credit to complete the purchase.
    \item In \textbf{Proxy > HTTP history}, locate the GET and POST requests sent to \texttt{/api/checkout}.
    \item Compare the JSON structures in both requests. Notice that the GET response contains a \texttt{chosen\_discount} parameter that is absent from the POST request.
    \item Right-click the \texttt{POST /api/checkout} request and select \textbf{Send to Repeater}.
    \item Add the \texttt{chosen\_discount} parameter to the JSON request body:
          \begin{lstlisting}[language=JSON]
{
    "chosen_discount": {
        "percentage": 0
    },
    "chosen_products": [
        {
            "product_id": "1",
            "quantity": 1
        }
    ]
}
    \end{lstlisting}
    \item Send the request and observe that adding the parameter does not cause an error.
    \item Change the \texttt{chosen\_discount} value to the invalid string \texttt{"x"}:
          \begin{lstlisting}[language=JSON]
{
    "chosen_discount": {
        "percentage": "x"
    },
    "chosen_products": [
        {
            "product_id": "1",
            "quantity": 1
        }
    ]
}
    \end{lstlisting}
          Send the request and observe the error indicating that the percentage must be a number. This suggests that the parameter is being processed by the application.
          \newpage
    \item Change the discount percentage to \texttt{100}:
          \begin{lstlisting}[language=JSON]
{
    "chosen_discount": {
        "percentage": 100
    },
    "chosen_products": [
        {
            "product_id": "1",
            "quantity": 1
        }
    ]
}
    \end{lstlisting}
    \item Send the request to apply the discount and solve the lab.
\end{enumerate}

\newpage
\nsection{Lab: Exploiting Server-Side Parameter Pollution in a Query String}

\begin{enumerate}
    \item In Burp's browser, trigger a password reset for the \texttt{administrator} user.
    \item In \textbf{Proxy > HTTP history}, locate the \texttt{POST /forgot-password} request and the related JavaScript file:
          \begin{lstlisting}
/static/js/forgotPassword.js
    \end{lstlisting}
    \item Right-click the \texttt{POST /forgot-password} request and select \textbf{Send to Repeater}.
    \item Resend the request to confirm that the response is consistent.
    \item Change the \texttt{username} parameter to an invalid username, such as \texttt{administratorx}. Send the request and observe the \texttt{Invalid username} error.
    \item Attempt to inject an additional parameter using a URL-encoded ampersand (\texttt{\%26}):
          \begin{lstlisting}
username=administrator%26x=y
    \end{lstlisting}
          Send the request and observe the \texttt{Parameter is not supported} error, suggesting that the internal API interpreted \texttt{x=y} as a separate parameter.
    \item Attempt to truncate the server-side query string using a URL-encoded hash character (\texttt{\%23}):
          \begin{lstlisting}
username=administrator%23
    \end{lstlisting}
          Send the request and observe the \texttt{Field not specified} error, suggesting that the internal query contains an additional \texttt{field} parameter.
    \item Inject a \texttt{field} parameter with an invalid value and truncate the remaining query string:
          \begin{lstlisting}
username=administrator%26field=x%23
    \end{lstlisting}
          Send the request and observe the \texttt{Invalid field} error, indicating that the injected parameter is recognized by the internal API.
    \item Brute-force valid values for the \texttt{field} parameter:
          \begin{enumerate}
              \item Right-click the request and select \textbf{Send to Intruder}.
              \item Set the payload position around the field value:
                    \begin{lstlisting}
username=administrator%26field=<PAYLOAD>%23
        \end{lstlisting}
              \item Under \textbf{Payload configuration}, select \textbf{Add from list}.
              \item Choose the built-in \textbf{Server-side variable names} payload list.
              \item Start the attack and review the results. Notice that the \texttt{username} and \texttt{email} payloads return HTTP \texttt{200} responses.
          \end{enumerate}
          \newpage
    \item Set the \texttt{field} parameter to \texttt{email}:
          \begin{lstlisting}
username=administrator%26field=email%23
    \end{lstlisting}
          Send the request and observe that it returns the original response, confirming that \texttt{email} is a valid field.
    \item In \textbf{Proxy > HTTP history}, inspect the \texttt{/static/js/forgotPassword.js} file. Identify the password reset endpoint:
          \begin{lstlisting}[language=JavaScript]
/forgot-password?reset_token=${resetToken}
    \end{lstlisting}
    \item Change the \texttt{field} parameter from \texttt{email} to \texttt{reset\_token}:
          \begin{lstlisting}
username=administrator%26field=reset_token%23
    \end{lstlisting}
          Send the request and retrieve the administrator's password reset token.
    \item In Burp's browser, navigate to the password reset endpoint and supply the retrieved token:
          \begin{lstlisting}
/forgot-password?reset_token=YOUR_TOKEN
    \end{lstlisting}
    \item Set a new password for the administrator account.
    \item Log in as \texttt{administrator} using the new password.
    \item Navigate to the \textbf{Admin panel} and delete \texttt{carlos} to solve the lab.
\end{enumerate}

\end{document}